\section{2024-11-20 U-I-Kennlinie von LEDs}

Ermittlung der Planckschen Konstante $h$

Erstellen Sie die U-I-Kennlinie der 4 LEDs. Verwenden Sie zur Messung die
Cassy-App.

Ermitteln und markieren Sie die Grenzspannung $U_G$.

Berechnen Sie die Energie $E=e\cdot U_G$ und die Frequenz des Lichtes
$f=\frac{c}{\lambda \text{max}}$ mit $c$ Lichtgeschwindigkeit und
$\lambda \text{max}$ Wellenlänge mit höchster Intensität.

Zeichnen Sie das $f$-$E$-Diagramm und ermitteln Sie gemäß $E=h\cdot f$ das
Planck'sche Wrikungsquantu $h$

% Load the Markdown module.
% \startluacode
% local kpse = require("kpse")
% kpse.set_program_name("luatex")
% \stopluacode
% \usemodule[t][markdown]
%
%  \startmarkdown
%
%  This is a paragraph of *Markdown text* with inline `\LaTeX`{=tex} code.
%
%  ``` {=tex}
%  This is a paragraph of \LaTeX{} code with inline \markinline{*Markdown text*}.
%  ```
%
%  \stopmarkdown
